% concatenated from Torsion_divisors_on_Fermat_varieties.tex
\documentclass[12pt,reqno]{article}
\usepackage{amsgen, amsmath, amsfonts, amsthm, amssymb, enumerate,amscd,graphics,dsfont,comment,tikz-cd}
\usetikzlibrary{intersections}

\oddsidemargin -.015in \evensidemargin -.015in \textwidth 6.6in \topmargin -.15in \textheight 9in
\newtheorem{defn}{Definition}[section]
\newtheorem{prop}[defn]{Proposition}
\newtheorem{lem}[defn]{Lemma}
\newtheorem{thm}[defn]{Theorem}
\newtheorem{cor}[defn]{Corollary}
\newtheorem{rem}[defn]{Remark}
\newtheorem {conj}[defn]{Conjecture}


\newcommand {\BL}{{\mathds L}}
\newcommand {\cf}{{\it cfr.}}
\newcommand {\ZZ}{{\mathds Z}}
\newcommand {\W}{{\bf W}}
\newcommand {\WW}{{\mathcal W}}
\newcommand {\KK}{{\mathcal K}}
\newcommand{\NN}{{\mathcal N}}
\newcommand {\B}{{\mathcal B}}
\newcommand {\X}{{\bf X}}
\newcommand {\XX}{{\mathcal X}}
\newcommand {\K}{{\mathds K}}
\newcommand {\E}{{\mathcal E}}
\newcommand {\C}{{\mathds C}}
\newcommand {\TR}{{\mathfrak R}}
\newcommand {\leg}{\Genfrac{(}{)}{}{}}
\newcommand {\A}{{\mathcal A}}
\newcommand {\F}{{\mathds F}}
\newcommand {\G}{{\mathds  G}}
\newcommand {\Z}{{\mathcal Z}}
\newcommand {\CC}{{\mathcal C}}
\newcommand {\Q}{{\mathds Q}}
\newcommand {\QQ}{{\bf Q}}
\newcommand {\R}{{\mathds R}}
\newcommand {\OO}{{\mathcal O}}
\newcommand {\lprime}{{\ell^2}}
\newcommand {\e}{{\mathbf e}}
\newcommand {\HH}{{\mathds H}}
\newcommand {\UH}{{\mathcal H}}
\newcommand {\LL}{{\mathcal L}}
\newcommand {\Y}{{\bf Y}}
\newcommand {\M}{{\mathcal M}}
\newcommand {\m}{{\mathfrak m}}
\newcommand {\V}{{\bf  V}}
\newcommand {\XY}{{\mathcal X}}
\newcommand {\PP}{{\mathfrak p}}
\newcommand {\ba}{{\mathfrak a}}
\newcommand {\YY}{{\mathcal Y}}
\newcommand {\BS}{{\bf S}_{0}}
\newcommand {\CS}{{\bf S}}
\newcommand {\EE}{{\bf E}}
\newcommand {\CP}{{\mathds P}} 
\newcommand {\TT}{{\mathcal T}}
\newcommand {\KI}{{K_{\infty}}}
\newcommand {\I}{{\mathcal I}}
\newcommand {\CI}{{{\mathds C}_{\infty}}}
\newcommand {\OI}{{{\mathcal O}_{\infty}}}
\newcommand {\D}{{\mathcal D}}
\newcommand {\QL}{{{\mathds Q}_{\ell}}}
\newcommand {\CPP}{{\mathds P} \times {\mathds P}} 
\newcommand {\GG}{{\mathcal G}}
\newcommand {\AB}{{\mathds A}}
\newcommand {\TL}{\tilde{L}}

\newcommand {\TCC}{\tilde{\mathcal C}}
\newcommand {\mE}{{\bar{\mathcal E}}}
\newcommand {\mF}{{\mathcal F}}
\newcommand {\bD}{{\tilde{D}}}
\newcommand  {\bfD}{{\tilde{f}_{\bD}}}
\newcommand {\fD}{{f_{D}}}
\newcommand {\LD}{{L_{\Delta}}}
\newcommand {\TK}{{\tilde{K}}}
\newcommand {\TX}{{\tilde{\Xi}}}
\newcommand {\TD}{{\tilde{D}}}

\newcommand {\RE}{{\mathrm{E}}}

\newcommand {\TQ}{{\tilde{Q}}}
\newcommand {\HCP}{{\hat{\mathds P}}}
\newcommand{\w}{{\omega}}
\newcommand{\TC}{\tilde{C}}

\newcommand{\BC}{\bar{C}}

\newcommand {\BCC}{\bar{\mathcal C}}
\newcommand {\TQQ}{{\tilde{\mathcal Q}}}
\newcommand {\QD}{Q^{\Delta}}
\newcommand {\TFF}{{\tilde{\mathcal F}}}
\newcommand {\FF}{{\mathcal F}}
\newcommand {\AF}{{\mathds A}}
\newcommand {\TP}{{\tilde{P}}}
\newcommand {\TRR}{{\tilde{R}}}
\newcommand {\qt}{{q_{\tau}}}

\newcommand {\DPR}{{D_{P-R}}}
\newcommand {\DQR}{{D_{Q-R}}}
\newcommand {\ddt}{{\frac{\partial}{\partial t}}}
\newcommand {\ddtbar}{{\frac{{\partial}}{\partial {\overline t}}}}

\newcommand{\bs}{\tilde{S}}

\newcommand{\tiota}{\tilde{\iota}}


\newcommand{\hG}{\hat{G}}
\newcommand {\MMR}{{\mathcal MR}}

\def\Ind{\operatorname{Ind}}
\def\pic{\operatorname{\rho}}
\def\Pic{\operatorname{Pic}}
\def\Ker{\operatorname{Ker}}
\def\End{\operatorname{End}}
\def\Supp{\operatorname{Supp}}
\def\lcm{\operatorname{lcm}}
\def\Sym{\operatorname{Sym}}
\def\gcd{\operatorname{gcd}}
\def\div{\operatorname{div}}
\def\deg{\operatorname{deg}}
\def\mod{\operatorname{mod}}
\def\dim{\operatorname{dim}}
\def\Spec{\operatorname{Spec}}
\def\Im{\operatorname{Im}}
\def\ord{\operatorname{ord}}
\def\reg{\operatorname{reg}}
\def\Hom{\operatorname{Hom}}
\def\Ext{\operatorname{Ext}}
\title{Torsion cycles on Fermat varieties}
\author{Ramesh Sreekantan\\Indian Statistical Institute Bangalore}
%\email{rameshsreekantan@gmail.com/rsreekantan@isibang.ac.in}
%\address{Indian Statistical Institute, Bangalore}

%\classification{11G15,11G55,14K22,14C25,14G35,19E15}
%\keywords{Motivic cycles, modular forms, Abelian surfaces, rational curves}

\begin{document}
	\baselineskip=17pt
	\maketitle
	
	\begin{abstract}
		
		A theorem of Manin and Drinfeld \cite{madr},\cite{drin} states that any divisor of degree $0$ on the cusps of a modular curve is torsion in the Jacobian. An elegant proof of this result was provided by Elkik \cite{elki} using mixed Hodge theory. 
		
		Rohrlich \cite{rohr} proved a generalization of this to Fermat curves. In this note we reprove his results along the lines of the work of Elkik. We then use the same methods to generalize it to higher codimensional null homologous cycles as well as higher Chow cycles on Fermat varieties. 
		
		\vspace{\baselineskip}
		{\bf Mathematical Sciences Classification (2020):} 14C25, 19E15, 14C30.
		
		
		
	\end{abstract}
	
	
	
%	\tableofcontents
	
	\section{Introduction.}
	
	A well known theorem of Manin and Drinfeld  \cite{madr,drin} states that any divisor of degree $0$ supported on the cusps of a modular curve is torsion in the Jacobian. This was generalized to Fermat curves by   Rohrlich \cite{rohr}. He proved the following theorem 
	
	\begin{thm}[Rohrlich] Let $F_d=F_d^1$  be the Fermat curve of degree $d$. 
		%
		$$	X_0^d+X_1^d+X_2^d=0.$$
		%
		Let $S_d$ be the set of $3d$ `Fermat points' - which are points determined by the hyperplanes $X_i=0$. These are the `trivial' solutions to Fermats Last Theorem. Then any divisor of degree $0$ supported on $S_d$ is torsion is the Jacobian. That is, if $D=\sum_{P \in S_d} a_P P$ such that $\sum_{P \in S_d} a_P=0$ then there exists a function $f \in K(F_d)$ such that for some $k \in {\mathds N}$,
		%
		$$\div(f)=kD$$
	\end{thm}
	
	There is an elegant `pure thought' proof of the original Manin-Drinfeld theorem due to Deligne and Elkik \cite{elki}. The idea is that the Abel-Jacobi map can be described in terms of extensions of mixed Hodge structures. Further, if $X$ is a modular curve, $S$ the set of cusps and $Y=X- S$ the open curve, then the rational Hodge structure on $H^1(Y)$ is a direct sum of the Hodge structure of pure weight $1$ on $H^1(X)$ and a Hodge structure of pure weight $2$. Once one knows this a short argument (Proposition \ref{splitting}) shows that any divisor of degree $0$ is torsion. 
	
	In this short paper  we show how the same idea can be applied to get a new proof of Rohrlich's theorem - though not quite as explicit as his is. However, one advantage is that this can be generalized to higher co-dimensions with only a little more work, though we get a weaker statement - only torsion in the intermediate Jacobian. 
	
	A linear subvariety of a Fermat hypersurface is a subvariety obtained by intersection with hyperplane sections. We have the following theorem (Theorem \ref{nullhomologoustheorem}, Corollary \ref{nullhomologouscorollary} ).
	
	
	  \begin{thm} Let $F^n_d$ be the Fermat hypersurface of degree $d$ and dimension $n$. Let $S_d^p=\cup_{|I|=p} C_I$ be the union of all the linear subvarieties of codimension $p$.  Then the exact sequence 
		%
		$$0 \rightarrow H^{2p-1}(F^n_d,\Q) \rightarrow H^{2p-1}(F^n_d-S_d^p,\Q) \rightarrow H^{2p}_{S_d^p}(F_d^n,\Q)^0 \rightarrow 0$$
		%
		is a split exact sequence of Hodge structures. Here 
		%
		$$H^{2p}_{S_d^p}(F_d^n,\Q)^0=\Ker\{H^{2p}_{S_d^p}(F^n_d,\Q) \rightarrow H^{2p}(F^n_d,\Q)\}$$
		%
		is the kernel of the cycle class map and is generated by linear combinations of the linear subvarieties  which are null-homologous. 
		
		As a result, any null-homologous cycle supported on $S_d^p$ is torsion in the $p^{th}$-intermediate Jacobian. 
	
	\end{thm}
	%
	A conjecture of Beilinson and Bloch asserts that the Abel-Jacobi map is injective for varieties defined over number fields, hence this would imply that the cycle is actually torsion in the Chow group itself. That is the case when $n=1$, which is Rohrlich's theorem. 
	
	 We then show  that the same methods can be used to obtain a statement about certain higher Chow cycles having torsion regulator (Theorem \ref{higherfermatsplitting}, Corollary \ref{higherchowtorsion})
	 
	 \begin{thm}
	 	Let $y$ be a higher Chow cycle  in the group $CH^p(F_d^n,q)_{\Q}=H^{2p-q}(F_d^n,\Q(p))$,
	 	%
	 	$$y=\sum (Z_I,f_{I_i},\dots, f_{I_q})$$
	 	%
	 	where $Z_I$ are linear subvarieties of codimension $p-q$ and $f_{I_j}$ are functions on $Z_I$ with  divisorial support on linear subvarieties of $Z_I$. Then the regulator $\reg(y)$ is $0$ in the Deligne cohomology group $H^{2p-q}_{\D}(F_d^n,\Q(p))$. 
	 
	 \end{thm}
	 
	 
	
	{\bf Acknowledgments:} I would like to thank the Korea Advanced Institute of Science and Technology (KAIST, Daejeon) for their hospitality when this work was started. I would like to acknowledge the use of AI - specifically Gemini 3 - for some of the arguments, literature and typographical and stylistic comments. I would like to thank Kannappan Sampath and Jinhyun Park for their comments. 
		
		\section{Hodge structures.}
		
		\subsection{Mixed Hodge structures.}
		
		According to Deligne \cite{deli}, a mixed Hodge structure (MHS) is an Abelian group $H_{\ZZ}$ with two filtrations. 
		
		\begin{itemize}
			\item  An increasing  {\em weight filtration} $W_{\bullet}$ on $H_{\Q}=H_{\ZZ} \otimes \Q$. 
			
			\item A decreasing {\em Hodge filtration} $F^{\bullet}$ on $H_{\C}=H_{\ZZ}\otimes \C$.
		\end{itemize}
	such that the graded pieces $Gr^{W_{\bullet}}_l = W_l/W_{l-1}$ 
    have a pure Hodge structure of weight $l$.  
		
		
	\subsection{Extensions of Mixed Hodge structures.}
	
	An extension of $B$ by $A$, where $A$ and $B$ are MHSs, is a short exact sequence of mixed Hodge structures  
	%
	$$e(B,A):0 \rightarrow A \rightarrow E \rightarrow B \rightarrow 0$$
	%
	$e(B,A)$ determines a class in the set $\Ext_{MHS}(B,A)=\Ext^1_{MHS}(B,A)$. The set $\Ext_{MHS}^1(B,A)$ is a group under the B\"ar sum. A class is trivial if and only if the exact sequence splits. 
	
	
	An extension is called a {\em separated} extension of MHS if the  highest weight of $A$ is less than the lowest weight of $B$. Carlson \cite{carl} showed that in this case that group can be identified with a generalized torus 
	
	\begin{thm}[Carlson] If $A$ and $B$ are  mixed Hodge structures such that the highest weight of $A$ is less than the lowest weight of $B$ - so any $e(B,A)$ in $\Ext(B,A)$ is separated. Then 
		%
		$$\Ext_{MHS}(B,A)=\Hom_{\C}(B,A)/(F^0 \Hom_{\C}(B,A)+\Hom_{\ZZ}(B,A))$$
		
		\end{thm}	 
		
		\subsection{Extensions and the Abel-Jacobi map.} 
	
	Suppose $X$ is a smooth projective variety and $D=\sum a_i D_i$ is a codimensional $p$ cycle. Let $|D|=\cup D_i$   denote the support of the cycle. Let $Y=X-|D|$. Then one has a long exact sequence of cohomology with support 
	%
	$$\dots \rightarrow  H^{2p-1}(X,\ZZ) \rightarrow H^{2p-1}(Y,\ZZ)\rightarrow H^{2p}_{|D|}(X,\ZZ)\rightarrow  H^{2p}(X,\ZZ)\rightarrow \dots$$
	%
	The term $H^{2p}_{|D|}(X,\ZZ)$ is pure of weight $2p$. Hence the cohomology group $H^{2p-1}$ has the weight filtration 
	\begin{itemize}
		\item $W_{2p-2}=\{0\} \subset W_{2p-1}=i^*(H^{2p-1}(X,\Q)\subset W_{2p}= H^{2p-1}(Y,\Q)$
		\item $Gr_{2p-1}^{W_{\bullet}}=W_{2p-1}=i^*(H^{2p-1}(X,\Q)$
		\item $Gr_{2p}^{W_{\bullet}}=H^{2p}_{|D|}(X,\Q)^0$, where 
		%
		$$H^{2p}_{|D|}(X,\Q)^0=\Ker\{H^{2p}_{|D|}(X,\Q) \rightarrow H^{2p}(X,\Q)\}$$
		
	\end{itemize}
	
	If $Z$ is a null-homologous cycle supported on $|D|$ then it determines a copy $\ZZ\cdot Z$ of $\ZZ(-p) \subset H^{2p}_{|D|}(X,\ZZ)$ which maps to $0$ in $H^{2p}(X,\ZZ)$ and  induces a short exact sequence 
	%
	$$0 \rightarrow H^{2p-1}(X,\ZZ) \rightarrow E_Z \rightarrow \ZZ(-p) \rightarrow 0$$
	%
	and $E_Z \in \Ext_{MHS}(\ZZ(-p),H^{2p-1}(X))$. 
	
	This is a separated extension of mixed Hodge structures. From Carlson's theorem we can see that $\Ext_{MHS}(\ZZ(-p),H^{2p-1}(X))$ is the same as the {\bf $p^{th}$ intermediate Jacobian.} 
	
	The {\bf (higher) Abel-Jacobi map} is given by the association 
	%
	$$Z \rightarrow [E_Z] \in \Ext_{MHS}(\ZZ(-p),H^{2p-1}(X)).$$
	%
	In particular, if $p=1$ this is the usual Abel-Jacobi map to the Jacobian of $X$. Hence we have 
	
	\begin{lem}
	A divisor $Z$ of degree $0$ is the divisor of a function if and only if the exact sequence
	%
	$$0 \rightarrow H^1(X) \rightarrow E_Z \rightarrow \ZZ(-1) \rightarrow 0$$
	%
	splits as an extension of MHSs. More generally it is torsion in the Jacobian  if and only if the sequence of $\Q$-MHSs splits.  
	
	\end{lem}
	
	
	This argument can be generalized to the following 
	
	\begin{prop} Let $X$ be a smooth projective variety and $|D|=\cup D_i$ a union of codimensional $p$ subvarieties. Let $Y=X-|D|$. Let 
		%
		$$H^{2p}_{|D|}(X,\Q)^0= \Ker\{H^{2p}_{|D|}(X,\Q) \rightarrow H^{2p}(X,\Q)\}$$
		%
		Then any null-homologous divisor $Z$ supported on $D$ is torsion in the $p^{th}$ intermediate Jacobian if and only if the exact sequence 
		%
		$$0 \rightarrow H^{2p-1}(X,\Q) \rightarrow  H^{2p-1}(Y,\Q) \rightarrow H^{2p}_{|D|}(X,\Q)^0 \rightarrow 0$$
	
		splits.  
		\label{splitting}
		
	\end{prop}
	\begin{proof} If $Z$ is a null-homologous cycle supported on $|D|$ then it determines a copy of $\Q(-p)$ in $H^{2p}_{|D|}(X,\Q)$. If the exact sequence is split then the same splitting implies that the class of $E_Z$ is split. Conversely, $H^{2p}_{|D|}(X,\Q)^0$ is generated by $\Q\cdot (D_i-D_j)|i \neq j$ so if $E_{D_i-D_j}$ is trivial for all $i,j$ then the sequence 
		%
			$$0 \rightarrow H^{2p-1}(X,\Q) \rightarrow  H^{2p-1}(Y,\Q) \rightarrow H^{2p}_{|D|}(X,\Q)^0 \rightarrow 0$$
		%
		is split. 
		\end{proof}
		
	For higher codimensional varieties it is not necessarily the case that a cycle which is trivial in the intermediate Jacobian is trivial in the Chow group. Hence from the splitting of the exact sequence one can only conclude that the cycle is $0$ in the intermediate Jacobian or is torsion if the $\Q$-mixed Hodge structure splits. A conjecture of Bloch \cite{blochinjectivity} and Beilinson asserts that if the variety is defined over a number field, such a cycle is necessarily torsion in the Chow group as well.  
		
		
		
		
		
		\section{Fermat varieties.}
		
		\subsection{Rohrlich's theorem and its generalizations.}
		
	From Proposition \ref{splitting}, 	in order to prove  Rohrlich's theorem we have to prove that the sequence 
	%
	$$0 \rightarrow H^1(F_d,\Q) \rightarrow H^1(F_d-S_d,\Q) \rightarrow H^{2}_{|S_d|}(F_d,\Q)^0 \otimes \Q \rightarrow 0$$ 
	%
	is split as a sequence of mixed Hodge structures. Equivalently we need to find a complement to the image of $H^1(F_d,\Q)$ in $H^1(F_d-S_d,\Q)$ which is a pure Hodge structure of weight $2$. 
	

	We first need some definitions. Let $F^n_d$ be the {\bf Fermat hypersurface of dimension $n$ and degree $d$} in $\CP^{n+1}$:
		%
		$$F^n_d:X_0^d+X_1^d+\dots+X_{n+1}^d=0$$
		%
		 
  For $I=\{i_1,\dots,i_p \} \subset \{0,\dots,n+1\}$ let $C_I=H_{i_1} \cap H_{i_2} \dots \cap H_{i_p} \cap F^n_d$, where $H_k$ is the hyperplane $X_k=0$. $C_I$ is a subvariety isomorphic to the Fermat variety $F_d^{n-p}$ and lies in $CH^p(F_d^n)$. We call an irreducible component of  such a subvariety a {\bf linear subvariety of codimension $p$}. If $n-p>1$ the linear subvarieties are irreducible so its only the $n=p$ case in which there could be components. However, as we also assume that $p=\frac{n+1}{2}$ this leaves us with only the case when $n=1$, namely the case of Rohrlich's theorem that is different. 
  
  We have the following theorem 
  
  \begin{thm} Let $F^n_d$ be the Fermat hypersurface of dimension $n$ and degree $d$. Let $S_d^p=\cup_{|I|=p} C_I$ be the union of all the linear subvarieties of codimension $p$.  Then the exact sequence 
	%
	$$0 \rightarrow H^{2p-1}(F^n_d,\Q) \rightarrow H^{2p-1}(F^n_d-S_d^p,\Q) \rightarrow H^{2p}_{S_d^p}(F_d^n,\Q)^0 \rightarrow 0$$
	%
	is a split exact sequence of Hodge structures. Here 
	%
	$$H^{2p}_{S_d^p}(F_d^n,\Q)^0=\Ker\{H^{2p}_{S_d^p}(F^n_d,\Q) \rightarrow H^{2p}(F^n_d,\Q)\}$$
	%
	is the kernel of the cycle class map and is generated by linear combinations of the linear subvarieties  which are null-homologous.
	\label{nullhomologoustheorem} 
	\end{thm}
	
	We get the following corollary: 
%	
 \begin{cor}
 	\label{nullhomologouscorollary}
	Any null homologous cycle supported on the set of linear subvarieties  is torsion in the intermediate Jacobian. 
\end{cor}
%
In particular, when $n=1$ and $p=1$ this is Rohrlich's theorem. Note that if $n>1$ the linear subvarieties are all irreducible but in the case of curves they are not. 
			

In order to prove the splitting we have to understand the cohomology of $F_d^n$. For that we follow Shioda \cite{shio}, who proved the Hodge conjecture for such varieties under some conditions. 

Let $\mu_d$ denote the group of $d^{th}$ roots of unity.  The group $\mu_d^{n+2}$ acts on $F_d^n$ by 
%
$$(\zeta_0,\dots,\zeta_{n+1})\cdot [X_0:\dots:X_{n+1}]=[\zeta_0X_0:\dots:\zeta_{n+1}X_{n+1}]$$
%
The diagonal acts trivially. Let $G_d^n=\mu_d^{n+2}/(diag)\simeq \mu_d^{n+1}$. The dual group of $G_d^n$  can be identified with 
%
$$\hG^n_d=\{\chi=(a_0,\dots,a_{n+1})|a_i \in \ZZ/d\ZZ \text{ such that } \sum a_i=0 \mod d\}$$ 
$$\chi((\zeta_0,\dots,\zeta_{n+1}))=\prod \zeta_i^{a_i}.$$

The action of $G^n_d$ on the Fermat varieties induces an action on related cohomology groups. Since $G^n_d$ also preserves $S_d^p$ it also induces a group action on the cohomology group $H^{2p}_{S_d^p}(F_d^n,\ZZ)$.

\subsection{$G_d^n$-module structure of $H^n(F_d^n,\C)$.}


We want to understand the decomposition of the cohomology groups $H^i(F^n_d,\Q)$, $H^i(F^n_d-S_d^p,\Q)$ and $H^i_{S_d^p}(F_d^n,\Q)$ as $G_d^n$ representations.


From the Lefschetz hyperplane theorem we know that 
%
$$H^i(F_d^n,\Q) = \begin{cases} \{0\} & i\neq n \text{ odd}\\
								\Q & i\neq n \text{ even} \end{cases}$$
								%
Hence the only `interesting' cohomology is when $i=n$. 

Recall that the primitive cohomology is the orthogonal complement of the hyperplane class. From above, all the primitive cohomology groups are trivial for $i\neq n$. When $n$ is odd, $H^n_{prim}(F_d^n,\Q)\simeq H^n(F_d^n,\Q)$ but when $n$ is even the primitive cohomology is of codimension $1$. 


For a character $\chi=\chi(a_0,a_1,\dots,a_{n+1})$, let $V(\chi)$ denote corresponding  $1$-dimensional representation.  Shioda has the following theorem 
 
\begin{thm}[Shioda] Let $A_d^n \subset \hG_d^n$ be the subset 
	%
	$$A_d^n=\{\chi=\chi(a_0,\dots,a_{n+1})|\, a_i \neq 0 \text{ for all } i\}$$
	%
	we have 	
 $$H^n_{prim}(F_d^n,\C)\simeq \bigoplus_{\chi \in A^n_d} V(\chi)$$ 		
	
\label{shiodatheorem}
	\end{thm}
	
	
	Null-homologous cycles of codimension $p$  determine an extension in 
	%
	$$\Ext^1_{MHS}(\Q(-p), H^{2p-1}(F_d^n,\Q)).$$
	%
	If $2p-1 \neq n$ then the cohomology group is trivial. So we can restrict ourselves to the case when $n=2p-1$ is odd  and in this case $H^n_{prim}=H^n$. In the section on motivic cycles we will need to consider the case when $n$ is even. 
	
	
	
	
This describes the action on the weight $(2p-1)$ part of the Hodge structure of $H^{2p-1}(F_d^n-S_d^p,\Q)$. 

\subsection{$G_d^n$-module structure of $H^{2p}_{S_d^p}(F_d^n,\Q)$.}

To understand the $Gr_{2p}^{W_{\bullet}}$ part of the Hodge structure we have to understand the $G_d^n$ action on $H^{2p}_{S_d^p}(F_d^n,\Q)$. This is the free $\Q$-vector space generated by the classes of $C_I$ in $H^{2p}(F_d^n,\Q)$, where $I=\{i_0,\dots,i_{p-1}\}$ is a subset of $\{0,\dots,n+1\}$ of cardinality $|I|=p$. Recall that we can assume $p=(n+1)/2$. 

There are two cases:

\noindent Case 1: $n=1$. Here the linear subvarieties  are the cuspidal points. Let $P$ be a cuspidal point in $C_I$ for some $I$. Here $C_I=C_{\{i\}}$ for  $i\in \{0,1,2\}$ so we will just use $i$. 

For any $P$ in $C_i$, the subgroup  
%
$$G_P=1\times\dots \times \mu_d \times \dots 1$$
%
with $\mu_d$ in the $i$ coordinate fixes $P$. The orbit of $P$ under the group $G_d^1$ is $C_i$. Note that $G_P$ is the same for all $P$ in $C_i$,  so we denote it by $G_i$. 

We would like to understand the representation $V_I$ corresponding to this orbit. Since $G_P=G_i$ acts trivially on $P$, the representation is 
%
$$V_i=\Ind_{G_i}^{G_d^1} ({\mathds 1})$$
%
From  Frobenius reciprocity the representation decomposes into those characters which restrict to the trivial character of $G_i$ 
%
$$V_i=\oplus_{\chi|_{G_i}={\mathds 1}} V(\chi)$$
%
These are those characters where $a_i=0$.

Let $B_i$ be the set of characters such that $a_i=0$. 
%
$$H^{2}_{S_d^1}(F_d^1,\Q)=\bigoplus_{i=0}^2 \bigoplus_{\chi \in B_i } V(\chi)$$
%
Note that $B_i \cap A_d^n=\emptyset$ for all $i$. 

\vspace{\baselineskip}


\noindent Case 2:$n >1$. In this case the $C_I$ are irreducible. $G_d^n$ preserves $C_I$ for any $C_I$ so the action on the class of $C_I$ is trivial. Hence $G_d^n$ acts trivially on $H^{2p}_{S_d^p}(F_d^n,\Q)$.



\begin{comment}


Let  $H_I \subset G_d^n$ be the subgroup defined by  
%
$$H_I=\{(\zeta_1,\zeta_2,\dots,\zeta_{n+1})|\zeta_i=1 \text { for all } i\notin I\}$$
%
is isomorphic to $(\ZZ/d\ZZ)^{n+1-p}$ and acts trivially on  $C_I$.

Let $\OO_I$ be the orbit of $C_I$ under $G_d^n$ and $V_I$ the corresponding representation. We have 




%
$$V_I=\Ind_{H_I}^{G_d^n}({\mathds 1})$$
%
of the trivial representation of $H_I$. Frobenius reciprocity implies 
%
$$V_I \simeq \bigoplus_{\chi|_{H_I}={\mathds 1}} V(\chi)$$
%
This is the set of characters $\chi(a_0,\dots,a_{n+1})$ which restrict to the trivial character  on $H_I$. Since $\chi(a_0,\dots,a_{n+1})(\zeta_0,\dots,\zeta_{n+1})=\prod \zeta_i^{a_i}$ these are precisely those characters such that $a_i=0$ when $i\in I$. 


Let $\chi(I)$ be the set of such characters. We then have 
%
$$H^{2p}_{S_d^p}(F_d^n,\C) \simeq \bigoplus_{|I|=p} \bigoplus_{\chi \in \chi(I)} V(\chi)$$




\subsection{Splitting of the Hodge Structure}

The characters $\chi$ are defined over the cyclotomic field $K=\Q(e^{\frac{2\pi i}{d}})$.
Let $T(\chi)$ in $K[G_d^n]$ be the projector corresponding to a character $\chi$. For a subset $I \subset G_d^n$ let 
%
$$T(I)=\sum_{\chi|_{H_{I}}={\mathds 1}} T(\chi)$$
	%
This is the projector corresponding to the representation $V(I)$ which is defined over $\Q$ so in fact $T(I)\in \Q[G_d^n]$.

Let $T_p=\sum_{|I|=p} T(I)$.
%
The characters appearing in the decomposition of $H^{2p}(F_d^p,\C)$ are those with all $a_i \neq 0$ while the characters appearing in $V(I)$ have at least $p$ of the $a_i=0$. Hence they are disjoint sets.  

$T_p$ acts on $H^{2p-1}(F_d^n-S_d^p,\Q)$ and vanishes on $i^*(H^{2p-1}(F_d^n,\Q))$.  The image, therefore is a pure Hodge structure of weight $2p$ and is complementary to $W_{2p-1}$. Hence the Hodge structure on $H^{2p-1}(F_d^n-S_d^p,\Q)$ splits as a sum of two Hodge structures. 

This finishes the proof of the theorem. 

\end{comment}


\subsection{Splitting of the Hodge Structure.}

We can now complete the proof of Theorem \ref{nullhomologoustheorem}.

\begin{proof} We first deal with the case $n>1$. Consider the element 
	%
	$$T=\frac{1}{|G_d^n|} \sum_{g \in G_d^n} g$$
	%
acting on $H^{2p-1}(F^n_d-S_d^p,\Q)$. Since 
%
$$H^{2p-1}(F_d^n,\C)\simeq \bigoplus_{\chi \in A^n_d} V(\chi)$$
%
for  $v \in V(\chi)$, 
 %
 $$T(v)= \frac{1}{|G_d^n|}    \sum_{g \in G_d^n}  g\cdot v = \frac{1}{|G_d^n|}  \sum_{g \in G_d^n} \chi(g)v=0$$
 % 
 since $\chi$ is a non-trivial character. 
 
Define a map
%
$$ H^{2p}_{S_d^p}(F_d^n,\Q) \stackrel{\phi}{\longrightarrow} H^{2p}_{S_d^p}(F_d^n,\Q)$$
 %
$$ \phi(v_0)=T(v) \text{ where } v \in \partial^{-1}(v_0)$$
%
This is well defined as if $v, v'\in \partial^{-1}(v_0)$, $v-v' \in H^{2p-1}(F_d^n,\Q)$ so $T(v-v')=-0$. 
 
  
 
 
 Hence $H^n(F_d^n,\C) \subset \Ker(T)$. $T$ is also a morphism of Hodge structures as it is induced by an automorphism of $F_d^n$. Further, since $G_d^n$ acts trivially on  $H^{2p}_{S_d^p}(F_d^n,\Q)$, $T$ acts by multiplication by $|G_d^n|=d^{n+1}$. The image of $T$ is a pure Hodge substructure of weight $2p$ of $H^{2p-1}(F^n_d-S_d^p,\Q)$ and is isomorphic to $H^{2p}_{S_d^p}(F_d^n,\Q)$ and complementary to $H^{2p-1}(F^n_d,\Q)$. Hence it provides a splitting of the exact sequence. 
 
 
 When $n=1$ the argument is similar. Let 
 %
 $$ T=\sum_{|I|=p} \sum_{\chi \in B_I} p_{\chi}$$
 %
 where $p_{\chi}$ is the projector corresponding to the character $\chi$. 
 %
 Since $B_I \cap A_d^n$ are disjoint $T$ projects on to the weight two piece of the Hodge structure on $H^1(F_d^n-S_d^1,\C)$ 
 
 \end{proof}
 
 \begin{rem} The case when $n=2p-1$ is the only case where one expects potentially interesting cycles. For instance $CH^1_{hom}(F_d^1), CH^2_{hom}(F_d^3), CH^3_{hom}(F_d^5)$ and in fact there do exist interesting non-torsion cycles in these groups as we will see in Section \ref{nontorsion}. 



\end{rem}

\section{Motivic cycles.}
 

A motivic cycle, or equivalently a higher Chow cycle, in the group $CH^p(X,q)=H^{2p-q}_{\M}(X,\ZZ(p))$ is  a codimensional $p$ cycle on $\Delta^q X=X \times \Delta^q$, where $\Delta^q$ is the $q$-simplex 
%
$$\Delta^q=\{(t_0,\dots,t_q) \in \AB^{q+1}| \sum t_i=1\}$$ 
%
subject to certain conditions with respect to the faces \cite{blocACKT,scho}. 


In some cases these cycles  can be represented by 
%
$$\sum (Z_i,f_{i_1},f_{i_2},\dots f_{i_q})$$ 
%
where $Z_i$ are codimensional $p$ cycles on $X \times \AB^q$ and $f_{i_j}$ are functions on $Z_i$ subject to the vanishing of the tame symbol of the functions.  In general, such a cycle will determine a higher Chow cycle but it is not necessarily the case that all higher Chow cycles have such a presentation. However in what follows we will only consider cycles of this type. 


A special case is when $q=1$. A cycle in the group $CH^p(X,1)=H^{2p-1}_{\M}(X,\ZZ(p))$ is represented by a sum 
%
$$\sum (Z_i,f_i)$$
%
where $Z_i$ are subvarieties of codimensional $p-1$ and $f_i$ are functions on $C_i$ subject to the condition that $\sum \div(f_i)=0$.  We will primarily focus on this group though our arguments should work in general. 

\subsection{Fermat varieties.}

In the case when $X$ is a Fermat hypersurface we can use Corollary \ref{nullhomologouscorollary} to construct  motivic cycles.  

Let  $Z$ be a linear subvariety of the type $C_I$ for some $I \subset \{0,1,\dots,n+1\}$ as before. Let $\{J\}$ be the sets such that $I \subset J$ and  $|J|=|I|+1$. Any null homologous divisor on $C_I$ supported  on the $C_J$ is torsion in the Jacobian as if $\dim(C_I)>1$ then $H^1(F^n_d,\Q)=\{0\}$ and if $\dim(C_I)=1$ we can apply Rohrlich's theorem. Hence it is not very difficult to find functions. 

Using these functions we can construct higher Chow cycles. For instance, in the group $CH^p(X,1)$ a cycle is represented by a sum $\sum (Z_i,f_i)$ with $\sum \div(f_i)=0$. We can construct cycles as follows.


Let $I=\{i_1,\dots,i_{p-2}\}$ be a subset of $\{0,\dots,n+1\}$. Let $a,b,c$ be three elements of the complement of $I$ in $\{0,\dots,n+1\}$ and consider the linear subvarieties of the type  $C_{a}=C_{I \cup \{a\}}$  and $C_{ab}=C_{I \cup \{a,b\}}$. Note that the linear subvarieties $C_a$ and $C_b$ meet at $C_{ab}=C_{ba}$.  


On $C_a$ we have a function $f_{bc}$ where 
	%
	$$\div(f_{bc})=C_{ab}-C_{ac}$$
	%
	and similar functions on $C_b$ and $C_c$. 
	
Then 
%
$$(C_a,f_{bc})+(C_b,f_{ca})+(C_c,f_{ab})$$ 
%
is an element of $CH^{p}(F_d^n,1)$ as 
%
$$\div_{C_a}(f_{bc})+\div_{C_b}(f_{ca})+\div_{C_c}(f_{ab})$$
$$ C_{ab}-C_{ac}+C_{bc}-C_{ba}+C_{ca}-C_{cb}=0$$
%
More generally one can construct cycles in $CH^p(F_d^n,q)$ along the same lines. Ross \cite{ross} constructed cycles in the group $CH^2(F_d^1,2)$ using pairs of functions supported on the Fermat points. 

There are some trivial elements in the higher Chow groups of the form $(Z,a_1,\dots,a_q)$ where $Z$ is a linear subvariety and $a_i$ are constant functions - which trivially satisfy the co-cycle conditions. Cycles like these are called {\bf decomposable} cycles as they come from the product structure on higher Chow groups. We do not want to consider such cycles.  





\subsection{Extensions and motivic cycles.}

Analogously to the situation of null-homologous cycles one can associate an extension to a motivic cycle \cite{scholl}.


A conjecture of Beilinson asserts that if $X$ is a smooth projective variety then 
%
$$CH^p(X,q)_{\Q}=H^{2p-q}_{\M}(X,\Q(p))=\Ext_{\M}(\Q(-p),h^{2p-q-1}(X,\Q))$$
%
where $\M$ is the conjectured category of mixed motives. 

It is known that in the category of mixed Hodge structures, 
%
$$\Ext_{MHS}(\Q(-p),H^{2p-q-1}(X,\Q))=H^{2p-q}_{\D}(X,\Q(p)).$$
%
The term on the right is the Deligne cohomology group which is the analogue of the Jacobian and there is a map called the regulator map from motivic cohomology to Deligne cohomology analogous to the Abel-Jacobi map. 

A motivic cycle, therefore, determines an extension in the conjectural category of mixed motives as well as determines an extension in the category of mixed Hodge structures and the realization in the category of mixed Hodge structures is the regulator map. 

To see how one associates an extension to a motivic cycle we follow Scholl \cite{scholl}, Section 3.1. Let $\Sigma^qX$ be the strict augmented simplicial scheme over $X$. 
%
$$\Delta^q X \leftarrow \bigsqcup^{[0,q]} \Delta^{q-1}X \mathrel{\substack{\leftarrow \\ \leftarrow}} \bigsqcup^{[1,q]} \cdots  \mathrel{\substack{\leftarrow \\ \cdot \\ \cdot  \\ \leftarrow}} \bigsqcup^{[q-1,q]} \Delta^0 X   $$
%
Higher Chow groups or motivic cohomology groups can be defined using these schemes. For any suitable cohomology theory $h^*_{\_}(\_)$ that extends to simplicial schemes with support, one has an exact sequence 
%
$$0 \rightarrow h^{2p-1}(\Sigma^q X) \rightarrow h^{2p-1}(\Sigma^q X-|D|_{\bullet})\rightarrow h^{2p}_{|D|_{\bullet}}(\Sigma^q X) \rightarrow h^{2p}(\Sigma^q X) $$
%
where for a higher Chow cycle $y$, $|D|$ denotes its support and $|D|_{\bullet}$ the closed simplicial subscheme of $\Sigma^q X$ determined by $|D|$. 

Scholl shows that using the tensor category of mixed realizations $\MMR$ of Deligne and Jannsen \cite{jann} it is possible to define mixed realizations of simplicial schemes with smooth components with supports and hence we can assume that $h^*_{\_}(\_)$ is the realization in the category of mixed realizations. One further has  \cite{scholl}, Section 3.1,
%
$$h^{*}(\Sigma^q X) \simeq h^{*-q}(X)$$ 
%
The cycle $y$ determines a copy of $\Q(-p)$ in $h^{2p}_{|D|_{\bullet}}(\Sigma^q X)$ which maps to $0$ in $h^{2p}(\Sigma^q X)$. So a higher Chow cycle determines a short exact sequence 
%
$$0 \rightarrow h^{2p-1-q}(X) \rightarrow E_y \rightarrow \Q(-p) \rightarrow 0$$
%
and in particular a class in $\Ext^1_{\MMR}(\Q(-p),h^{2p-1-q}(X))$. 


We now have an analogue of Proposition \ref{splitting}

	\begin{prop} Let $X$ be a smooth projective variety and $|D|=\cup D_i$ a union of codimensional $p$ subvarieties of $\Delta^q(X)$. Let $|D|_{\bullet}$ be the simplicial subscheme of $\Sigma^q X$ induced by $|D|$ 
	%
	$$H^{2p}_{|D|_{\bullet}}(\Sigma^q X,\Q)^0= \Ker\{H^{2p}_{|D|_{\bullet}}(\Sigma^q X,\Q) \rightarrow H^{2p}(\Sigma^q X,\Q)\}$$
	%
	Then any higher cycle  $Z$ in $Z^p(X,q)$ supported on $|D|$ is torsion in the Deligne cohomology $H^{2p-q}_{\D}(X,\Q(p)$  if the exact sequence 
	%
	$$0 \rightarrow H^{2p-1}(\Sigma^q X,\Q) \rightarrow  H^{2p-1}(\Sigma^q X - Y_{\bullet},\Q) \rightarrow H^{2p}_{|D|}(\Sigma^q X,\Q)^0 \rightarrow 0$$
	%
	splits.  
	
	\label{highersplitting} 
	
\end{prop}
\begin{proof} The proof is similar to the earlier case. If $y$ is a higher Chow cycle in the
	group $CH^p(X,q)=H^{2p-q}_{\M}(X,\Q(p))$ supported on $|D|$ then it determines a copy 
	 of $\Q(-p)$ in $H^{2n}_{|D|}(\Sigma^q X,\Q)$ and its pull-back determines an extension 
	 $E_y$ in $\Ext(\Q(-p),H^{2p-q-1}(X))$. If the exact sequence is split then the same 
	 splitting implies that the class $E_y$ is split. 
\end{proof}
  

\subsection{Motivic cycles on Fermat Varieties.}

We now show that any motivic cycle supported on the linear hypersurfaces of a Fermat variety have regulator which is torsion in the Deligne cohomology. This is the higher Chow version of Rohrlich's theorem. The idea, as before, is to show that in this case the exact sequence in Proposition \ref{highersplitting} splits. 

We consider cycles in the group $CH^p(F_d^n,q)$ with $p-q \geq 0$ and $n>1$. The cycles $y$ are of the following type:
%
$$y=\sum (Z_I,f_{I_1},\dots, f_{I_q})$$
%
where $Z_I$ are linear subvarieties of codimension $p-q$ corresponding to a subset $I$ of $\{0,\dots,n+1\}$ and $f_{I_k}$ are functions on the $Z_I$ such that their divisors are supported on cycles of the form $Z_{I\cup \{j\}}$ for some $j\in \{0, \dots,n+1\}, j \notin I$ and further they satisfy a co-cycle condition ensuring that the sum is in the motivic cohomology group. 

We want to discount the decomposable cycles coming from linear subvarieties. In order to do that we work with primitive cohomology. 


The group $G_d^n$ stabilizes $Z_I$ and the support of the divisors of $f_{I_j}$ as they are all linear subvarieties. Hence similar to before since $n>1$ the group acts trivially on the classes of a higher cycle  $y=\sum_I (Z_I,f_{I_i},\dots,f_{I_q})$ in  $h^{2p}_{|D|_{\bullet}}(\Sigma^q X)$, where $|D|_{\bullet}$ is the simplicial subscheme determined by cycles of the type above. 

\begin{thm} Let $F^n_d$ be the Fermat hypersurface of dimension $n$.
	%
	$$F^n_d:X_0^d+X_1^d+\dots+X_{n+1}^d=0$$
	%
	For $I=\{i_1,\dots,i_p \} \subset \{0,\dots,n+1\}$ let $C_I=H_{i_1} \cap H_{i_1} \dots \cap H_{i_p} \cap F^n_d$, where $H_k$ is the hyperplane $X_k=0$. $C_I$ is a subvariety isomorphic to the Fermat variety $F_d^{n-p}$. Let $S_d^p=\bigcup_{|I|=p} C_I$.
	
	
	Then the exact sequence 
	%
	$$0 \rightarrow H^{2p-1}_{prim}(\Sigma^q F^n_d,\Q) \rightarrow H^{2p-1}_{prim}(\Sigma^q F^n_d-S_d^p,\Q) \rightarrow H^{2p}_{prim,S_d^p}(F_d^n,\Q)^0 \rightarrow 0$$
	%
	is a split exact sequence of Hodge structures. Here 
	%
	$$H^{2p}_{prim,S_d^p}(F_d^n,\Q)^0=\Ker\{H^{2p}_{S_d^p}(F^n_d,\Q) \rightarrow H^{2p}_{prim}(F^n_d,\Q)\}$$
	%
	is the kernel of the cycle class map and is generated by linear combinations of the $C_I$ which are null-homologous.
	\label{higherfermatsplitting} 
\end{thm}

\begin{proof} The proof is identical to Theorem \ref {nullhomologoustheorem} as one can once again consider the operator $T=\sum_{g \in G_d^n} g$ and for exactly the same reasons it produces a splitting. 
	
\end{proof}
We get the following corollary: 
%	
\begin{cor}
	Let $y$ be a motivic cycle in $H^{2p-q}(F_d^n,\Q(p)=CH^p(F_d^n,q)_{\Q}$ of the form 
	% 
	$$y=\sum (Z_I,f_{I_i},\dots, f_{I_q})$$
	%
 where $Z_I$ are linear subvarieties and $f_{I_j}$ functions on $Z_I$ with divisorial support on linear subvarieties of $Z_I$. Then $\reg(y)$ is $0$  in the Deligne cohomology group $H^{2p-q}_{\D}(F_d^n,\Q(p))$. 
	\label{higherchowtorsion}
\end{cor}

\begin{rem} The Bloch-Beilinson conjectures on injectivity of the regulator map  imply that these cycles are torsion in the Chow group itself. 
	\end{rem}
	
	\begin{rem} The argument does not work when $n=1$, that is Fermat curves, as in that case the divisor of a function is supported on points which is not necessarily fixed by the action of $G_d^1$ since the $0$-dimensional linear subvarieties are not irreducible. 
		 In fact there are examples of non-trivial cycles. In the case of  $CH^2(F_d^n,2)$ - Ross \cite{ross}  constructed cycles in the group $CH^2(F_d^n,2)$ which are non-trivial using functions with divisors on the Fermat points. 
		\end{rem} 
		
		
\begin{rem} Again, like in the case of null-homologous cycles, the only really interesting higher Chow groups are $CH^p(F_d^{2p-q-1},q)$,  for example $CH^2(F_d^2,1),\, CH^2(F_d^1,2),\, CH^3(F_d^4,1), \dots$. These are the cases where the Deligne cohomology is determined by the middle cohomology group. 
	
Once again is Section \ref{nontorsion} we will see that there are possibly non-torsion cycles in these groups coming from the product structure. 
	\end{rem}



\section{Non-torsion cycles on  Fermat varieties.}
\label{nontorsion}

The above theorems show that certain `natural' constructions do not produce interesting cycles on Fermat varieties. This does not mean that the Chow groups of Fermat varieties are trivial. 

There is an inductive structure on Fermat Varieties which has interesting implications for  cycles on Fermat varieties. There is a rational map 
%
$$F_d^r \times F_d^s \dashrightarrow F^{r+s}_d$$
%
defined by 
%
$$[x_0:\dots:x_{r+1}],[y_0:\dots:y_{s+1}] \dashrightarrow [x_0 y_{s+1}:\dots:x_r y_{s+1}:\epsilon x_{r+1}y_0:\dots:\epsilon x_{r+1}y_s]$$ 
%
where $\epsilon^{d}=-1$. As a result we have a map 
%
$$F_d^1 \times F_d^1 \times F_d^1 \dashrightarrow F_d^3$$

In the group $CH^2_{hom}((F_d^1)^3)_{\Q}$ there is a well known cycle called the modified diagonal cycle of Gross and Schoen \cite{grsc}. If $e$ is a point on $F_d^1$, 
%
$$\Delta_e=\Delta-\Delta_{12}-\Delta_{13}-\Delta_{23}+\Delta_1+\Delta_2+\Delta_3$$
%
where $\Delta_{ij}$ and $\Delta_i$ is the diagonal with $e$ in the complementary positions.
This is called the Gross-Schoen modified diagonal cycles. It is well known that this cycle is in general non-trivial in the Chow group of null-homologous cycles on $(F_d^1)^3$ \cite{otsu}. 

We can consider the image of this cycle under the map induced by the inductive structure to get a null-homologous cycle in $CH^2_{hom}(F_d^3)$ which is likely to be non-torsion. More generally there are modified diagonal cycles in $CH^p_{hom}((F_d^1)^{2p-1})$ studied by Gross-Schoen \cite{grsc} and more recently by Eskandari \cite{eska}. It is likely that these cycles will also map to non-torsion null homologous cycles on $F_d^{2p-1}$ 

Similarly, in the group $CH^2(F^1_d \times F^1_d,1)$, Sarkar \cite{sark} constructed higher Chow cycles using functions on the Fermat curve supported on the cusps and a construction of Bloch's \cite{blocirvine}.

As we mentioned earlier, Ross \cite{ross} constructed non-torsion cycles in the groups $CH^2(F_d^1,2)$ using functions supported on the cusps. 
















\bibliographystyle{alpha}
\bibliography{Fermatvarieties.bib}








		
\end{document}
